\appendix
    \chapter{Derivation of the Hamiltonian for the chain of Rydberg atoms in optical tweezers}
    \label{appendix:derivation_Hamiltonian}
As described in the main text, the assumption of small displacements $\{\delta\mathbf{R}_l\}$ compared to the distance between the tweezers motivates the Taylor expansion of the interaction potential to second order.
% By redefining the trap center the first order term can be absorbed. Consider a function $f(r)$ which is expanded up to first order in $r$, i.e. $f(r+\delta r)\approx f(r)+ \partial_rf(r)\cdot\delta r$, if we now add a trap-like potiential
% \begin{align*}
%     \frac12 m\omega^2 \,\delta r^2 + f(r) +\partial_rf(r)\,\delta r&= \frac12 m\omega^2\,(\delta r+\frac{\partial_rf(r)}{m\omega^2})^2+ \frac12 \,\frac{(\partial_rf(r))^2}{m\omega^2} + f(r)\\
%     &=\vcentcolon  \frac12 m\omega^2 \,\delta\tilde{r}^2 +\text{const.},
% \end{align*}
% Thus by discarding the first order the Hamiltonian up to second order reads
\begin{align*}
    \mathcal{H}_I(\{\mathbf{R}_l+\delta\mathbf{R}_l\})&=\mathcal{H}_I(\{\mathbf{R}_l\})+\sum_n\frac{\partial \mathcal{H}_I(\{\mathbf{R}_l\})}{\partial \mathbf{R}_n}\,\delta\mathbf{R}_n +  \frac12\sum_{n,m}\delta\mathbf{R}_n\,\frac{\partial^2 \mathcal{H}_I(\{\mathbf{R}_l\})}{\partial \mathbf{R}_n\partial\mathbf{R}_m}\,\delta\mathbf{R}_m\\
    &=\mathcal{H}_I(\{\mathbf{R}_l\})+\sum_n\frac{\partial \mathcal{H}_I(\{\mathbf{R}_l\})}{\partial \mathbf{R}_n}\,\delta\mathbf{R}_n +  \frac12\sum_{n\neq m}\delta\mathbf{R}_n\,\frac{\partial^2 \mathcal{H}_I(\{\mathbf{R}_l\})}{\partial \mathbf{R}_n\partial\mathbf{R}_m}\,\delta\mathbf{R}_m + \frac12\sum_{n= m}\delta\mathbf{R}_n\,\frac{\partial^2 \mathcal{H}_I(\{\mathbf{R}_l\})}{\partial \mathbf{R}_n\partial\mathbf{R}_m}\,\delta\mathbf{R}_m\\
    &=\mathcal{H}_I(\{\mathbf{R}_l\})+\sum_n\frac{\partial \mathcal{H}_I(\{\mathbf{R}_l\})}{\partial \mathbf{R}_n}\,\delta\mathbf{R}_n +  \frac12\sum_{n\neq m}\;\sum_{i,j;\;i> j}\delta\mathbf{R}_n\,\frac{\partial^2 V_{i,j}(\{\mathbf{R}_l\})}{\partial \mathbf{R}_n\partial\mathbf{R}_m}\,\delta\mathbf{R}_m  \\&+\frac12\sum_{n}\sum_{i,j;\;i> j}\delta\mathbf{R}_n\,\frac{\partial^2 V_{i,j}(\{\mathbf{R}_l\})}{\partial \mathbf{R}_n^2}\,\delta\mathbf{R}_n \\
    &=\mathcal{H}_I(\{\mathbf{R}_l\})+\sum_n\frac{\partial \mathcal{H}_I(\{\mathbf{R}_l\})}{\partial \mathbf{R}_n}\,\delta\mathbf{R}_n +  \frac12\sum_{n\neq m}\;\sum_{i,j;\;i> j}\delta\mathbf{R}_n\,\frac{\partial^2 V_{i,j}(\{\mathbf{R}_l\})}{\partial \mathbf{R}_n\partial\mathbf{R}_m}\,\delta\mathbf{R}_m \,(\delta_{n=i,m=j}+\delta_{n=j,m=i})  \\
    &+\frac12\sum_{n}\sum_{i,j;\;i> j}\delta\mathbf{R}_n\,\frac{\partial^2 V_{i,j}(\{\mathbf{R}_l\})}{\partial \mathbf{R}_n^2}\,\delta\mathbf{R}_n\,(\delta_{n=i}+\delta_{n=j}) \\
\end{align*}
In order to allow for a quantum mechanical description the displacements $\delta\mathbf{R}_l$ are treated as operators. The second order term can be simplified by exploiting the symmetry of the second derivative and the fact that $V_{i,j}$ only depends on $\mathbf{R}_i$ and $\mathbf{R}_j$. Thus
\begin{align*}
   \sum_{i,j;\;i> j}\left(\frac{\partial^2 V_{i,j}(\{\mathbf{R}_l\})}{\partial \mathbf{R}_i^2}\,\delta\mathbf{R}_i^2+\frac{\partial^2 V_{i,j}(\{\mathbf{R}_l\})}{\partial \mathbf{R}_j^2}\,\delta\mathbf{R}_j^2\right)&=\sum_{i,j;\;i\neq j}\frac{\partial^2 V_{i,j}(\{\mathbf{R}_l\})}{\partial \mathbf{R}_i^2}\,\delta\mathbf{R}_i^2
\end{align*}
and
\begin{align*}
    \sum_{i,j;\;i> j}\left(\delta\mathbf{R}_i\,\frac{\partial^2 V_{i,j}(\{\mathbf{R}_l\})}{\partial \mathbf{R}_i\partial\mathbf{R}_j}\,\delta\mathbf{R}_j + \delta\mathbf{R}_j\,\frac{\partial^2 V_{j,i}(\{\mathbf{R}_l\})}{\partial \mathbf{R}_j\partial\mathbf{R}_i}\,\delta\mathbf{R}_i\right)&= \sum_{i,j;\;i\neq j}\delta\mathbf{R}_i\,\frac{\partial^2 V_{i,j}(\{\mathbf{R}_l\})}{\partial \mathbf{R}_i\partial\mathbf{R}_j}\,\delta\mathbf{R}_j
\end{align*}
The interaction can therefore be written in the following form
\begin{align*}
    \mathcal{H}_I(\{\mathbf{R}_l+\delta\mathbf{R}_l\})&=\mathcal{H}_I(\{\mathbf{R}_l\})+\sum_i\frac{\partial \mathcal{H}_I(\{\mathbf{R}_l\})}{\partial \mathbf{R}_i}\,\delta\mathbf{R}_i + \frac12\, \sum_{i,j;\;i> j}\frac{\partial^2 V_{i,j}(\{\mathbf{R}_l\})}{\partial \mathbf{R}_i\partial\mathbf{R}_j}\, \left(\delta\mathbf{R}_i\,\delta\mathbf{R}_j+\delta\mathbf{R}_j\,\delta\mathbf{R}_i\right) \\
    &+\frac12\,\sum_{i,j;\;i\neq j}\frac{\partial^2 V_{i,j}(\{\mathbf{R}_l\})}{\partial \mathbf{R}_i^2}\,\delta\mathbf{R}_i^2\\
    &=\mathcal{H}_I(\{\mathbf{R}_l\})+\sum_i\frac{\partial \mathcal{H}_I(\{\mathbf{R}_l\})}{\partial \mathbf{R}_i}\,\delta\mathbf{R}_i + \frac12\, \sum_{i,j;\;i\neq j}\frac{\partial^2 V_{i,j}(\{\mathbf{R}_l\})}{\partial \mathbf{R}_i\partial\mathbf{R}_j}\,\delta\mathbf{R}_i\,\delta\mathbf{R}_j \\
    &+\frac12\,\sum_{i,j;\;i\neq j}\frac{\partial^2 V_{i,j}(\{\mathbf{R}_l\})}{\partial \mathbf{R}_i^2}\,\delta\mathbf{R}_i^2
\end{align*}
By introducing the matrix $U$ the expanded interaction Hamiltonian takes the compact form
\begin{align*}
    \mathcal{H}_I(\{\mathbf{R}_l+\delta\mathbf{R}_l\})&=\mathcal{H}_I(\{\mathbf{R}_l\})+\sum_n\frac{\partial \mathcal{H}_I(\{\mathbf{R}_l\})}{\partial \mathbf{R}_n}\,\delta\mathbf{R}_n + \sum_{n,m}U_{n,m}\,\delta\mathbf{R}_n\,\delta\mathbf{R}_m
\end{align*}
Where the constant energy shift $\mathcal{H}_I(\{\mathbf{R}_l\})$ can be ignored and the matrix $U$ is defined as
\begin{align*}
    U_{n,m}&=
        \frac12\,\frac{\partial^2 V_{n,m}(\{\mathbf{R}_l\})}{\partial \mathbf{R}_n\partial\mathbf{R}_m}\quad\text{if}\:n\neq m\\
    U_{n,n} &= \frac12\,\sum_{j;\,j\neq n}\frac{\partial^2 V_{n,j}(\{\mathbf{R}_l\})}{\partial \mathbf{R}_n^2}
\end{align*}
% Swithing again to the reduced writting, where $V_{ij}$ is only writen as a function of the variables of which it actually depends, reads
% \begin{align*}
%     U_{n,m}&=
%         \frac12\,\frac{\partial^2 V_{n,m}(\mathbf{R}_n,\mathbf{R}_m)}{\partial \mathbf{R}_n\partial\mathbf{R}_m}\quad\text{if}\:n\neq m\\
%     U_{n,n} &= \frac12\,\sum_{j;\,j\neq n}\frac{\partial^2 V_{n,j}(\mathbf{R}_n,\mathbf{R}_j)}{\partial \mathbf{R}_n^2}
% \end{align*}
For an interaction that only depends on the relative distance between the atoms $\mathbf{r}_{ij}=\mathbf{R}_i-\mathbf{R}_j$, i.e. $V_{i,j}(\{\mathbf{R}_l\})=V_{i,j}(\mathbf{R}_i-\mathbf{R}_j)$, the matrix $U$ can be further specified. Using the fact that
\begin{align*}
    \frac{\partial \mathbf{r}_{ij}}{\partial \mathbf{R}_i}=\mathbb{1}=-\frac{\partial \mathbf{r}_{ij}}{\partial \mathbf{R}_j}
\end{align*}
the derivatives in $U$ can be written as 
\begin{align*}
    \frac{\partial^2 V_{n,m}(\mathbf{r}_{nm})}{\partial \mathbf{R}_n\partial\mathbf{R}_m}&=\frac{\partial}{\partial  \mathbf{R}_n }\,\left( \frac{\partial V_{n,m}(\mathbf{r}_{nm})}{\partial\mathbf{r}_{nm}}\,  \frac{\partial \mathbf{r}_{nm}}{\partial\mathbf{R}_{m}}\right)\\
    &= \frac{\partial^2 V_{n,m}(\mathbf{r}_{nm})}{\partial\mathbf{r}_{nm}^2}\,  \frac{\partial \mathbf{r}_{nm}}{\partial\mathbf{R}_{m}}\frac{\partial \mathbf{r}_{nm}}{\partial\mathbf{R}_{n}}+\frac{\partial V_{n,m}(\mathbf{r}_{nm})}{\partial\mathbf{r}_{nm}}\,  \frac{\partial^2 \mathbf{r}_{nm}}{\partial\mathbf{R}_{n}\partial\mathbf{R}_{m}}\\
    &=-\frac{\partial^2 V_{n,m}(\mathbf{r}_{nm})}{\partial\mathbf{r}_{nm}^2}
\end{align*}
and
\begin{align*}
    \frac{\partial^2 V_{n,m}(\mathbf{r}_{nm})}{\partial \mathbf{R}_m^2m}
    &= \frac{\partial^2 V_{n,m}(\mathbf{r}_{nm})}{\partial\mathbf{r}_{nm}^2}\,  \frac{\partial \mathbf{r}_{nm}}{\partial\mathbf{R}_{m}}\frac{\partial \mathbf{r}_{nm}}{\partial\mathbf{R}_{m}}\\
    &=\frac{\partial^2 V_{n,m}(\mathbf{r}_{nm})}{\partial\mathbf{r}_{nm}^2}
\end{align*}
Thus
\begin{align*}
    U_{n,m}&=
        -\frac12\,\frac{\partial^2 V_{n,m}(\mathbf{r}_{nm})}{\partial\mathbf{r}_{nm}^2}\quad\text{if}\:n\neq m\\
    U_{n,n} &= \frac12\sum_{j;\,j\neq n}\frac{\partial^2 V_{n,j}(\mathbf{r}_{nj})}{\partial\mathbf{r}_{nj}^2}
\end{align*}
At this point I can eliminate the linear term of the Hamiltonian by introducing the trap potential 
\begin{align*}
    \mathcal{H}_0(\{\delta\mathbf{R}_l\}) = \sum_i\frac{\mathbf{P}_i^2}{2m}+ \frac12 m \omega^2 \,\delta\mathbf{R}_i^2,
\end{align*}
where $\mathbf{P}_i$ is the momentum of the $i$th atom. I can now shift the considered displacements from the trap centers by a constant $\{\delta\mathbf{R}_l\}=\{\delta\tilde{\mathbf{R}}_l+\mathbf{\Delta}_l\}$, where $\delta\mathbf{R}_l$ and $\delta\tilde{\mathbf{R}}_l$ are operators and $\mathbf{\Delta}_l$ are just constant vectors. Neglecting all constant terms the Hamiltonian minus the momentum term $\tilde{\mathcal{H}}=\mathcal{H}-\sum_i\mathbf{P}_i/2m$ up to second order reads
\begin{align*}
    \tilde{\mathcal{H}}(\{\delta\mathbf{R}_l\})&= \sum_i \frac12 m \omega^2 \,(\delta\tilde{\mathbf{R}}_i+\mathbf{\Delta}_i)^2+ \sum_i\frac{\partial \mathcal{H}_I(\{\mathbf{R}_l\})}{\partial \mathbf{R}_i}\,(\delta\tilde{\mathbf{R}}_i+\mathbf{\Delta}_i) + \sum_{i,j}(\delta\tilde{\mathbf{R}}_i+\mathbf{\Delta}_i)\,U_{i,j}\,(\delta\tilde{\mathbf{R}}_j+\mathbf{\Delta}_j)\\
    &= \sum_i \frac12 m \omega^2 \,\delta\tilde{\mathbf{R}}_i^2  + \sum_{i,j}\delta\tilde{\mathbf{R}}_i\,U_{i,j}\,\delta\tilde{\mathbf{R}}_j + \sum_i\left(\frac{\partial \mathcal{H}_I(\{\mathbf{R}_l\})}{\partial \mathbf{R}_i} +\sum_j\,\left( 2\, U_{j,i} + m\omega^2\,\delta_{ij}\right)\,\mathbf{\Delta}_j  \right)\,\delta\tilde{\mathbf{R}}_i,
\end{align*}
where it was used that $U$ is symmetric. As the matrix $\left( 2\, U_{j,i} + m\omega^2\,\delta_{ij}\right)$ is invertible, a set of vectors $\{\mathbf{\Delta}_n\}$ can be found such that the linear term vanishes.
% The linear term vanishes, if $\mathbf{\Delta}_n$ is chosen such that
% \begin{align*}
%     \mathbf{\Delta}_n = -\left( 2\,\sum_m\delta\tilde{\mathbf{R}}_m\, U_{m,n} + m\omega^2\right)^{-1} \,\frac{\partial \mathcal{H}_I(\{\mathbf{R}_l\})}{\partial \mathbf{R}_n}
% \end{align*}
Thus
\begin{align*}
    \mathcal{H}(\{\delta\tilde{\mathbf{R}}_l\})&= \sum_n\frac{\mathbf{P}_n^2}{2m}+ \sum_n\frac12 m\omega^2\,\delta\tilde{\mathbf{R}}_n^2+ \sum_{n,m}U_{n,m}\,\delta\tilde{\mathbf{R}}_n\,\delta\tilde{\mathbf{R}}_m
\end{align*}
At this point I introduce the quantized displacements in form oft the quadrature operators 
\begin{align*}
    \delta\tilde{\mathbf{R}}_n &= \sqrt{\frac{\hbar}{2m\omega}}\,(a_n + a_n^\dagger)\\
    \mathbf{P}_n &= i\,\sqrt{\frac{\hbar m\omega}{2}}\,(a_n^\dagger - a_n)
\end{align*}
% If we exploit that the interactions realized in $U$ are much smaller compared to the trap frequency, i.e. $U_{n,m}\ll m\omega^2$, the creation and annhiliation of phonons is supressed and the dynamics only involves phonon hopping, thus combinations $aa$ or $a^\dagger a^\dagger$ are discarded. 
With that the Hamiltonian gets its final form
\begin{align*}
    \mathcal{H} &= \sum_n \hbar\omega\,a_n^\dagger a_n +\frac{\hbar}{2m\omega}\, \sum_{n,m} U_{n,m}\,\left(a_n\,a_m + a_n^\dagger a_m + a_n\,a_m^\dagger + a_n^\dagger a_m^\dagger \right)
    % &=\vcentcolon\sum_n \hbar\omega\,a_n^\dagger a_n +\sum_{n,m} \tilde{U}_{n,m}\,a_n^\dagger a_m
\end{align*}\hfill\newpage

\chapter{Symmetry and irreducibility check}
    \label{appendix:symmetry}

In order to evaluate wether a matrix is irreducible one has check whether there are unitary matrices that specific matrix. It is helpful to decompose the matrices into tensor products of pauli matrices. for 4x4 matrices it is useful to know the commutator of the following tensor products
\begin{align*}
    \left[(\sigma_i\otimes\sigma_j), (\sigma_n\otimes\sigma_m)  \right]&= \half \,[\sigma_i,\sigma_n]\otimes\{\sigma_j,\sigma_m\}+\half\,\{\sigma_i,\sigma_n\}\otimes[\sigma_j,\sigma_m]
\end{align*}
in order to evaluate this expression, I distinguish three exclusive cases. First let at least one of the pauli matrices be the identity. Without loss of generality let $\sigma_i=\mathbb{1}$. Then it follows
\begin{align*}
    \left[(\sigma_i\otimes\sigma_j), (\sigma_n\otimes\sigma_m)  \right]&= \sigma_n\otimes[\sigma_j,\sigma_m]
\end{align*}
This vanishes if one matrix of the $j,m$ pair is the identity or if they are the same.
The other case is $\sigma_i=\sigma_n$ or $\sigma_j=\sigma_m$. Here one receives a similar result. Let $\sigma_i=\sigma_n$, then
\begin{align*}
    \left[(\sigma_i\otimes\sigma_j), (\sigma_n\otimes\sigma_m)  \right]&= \mathbb{1}\otimes[\sigma_j,\sigma_m]
\end{align*}
which equals to zero under the same premises as in case 1.
Finally if none of these two cases is true the result is
\begin{align*}
    \left[(\sigma_i\otimes\sigma_j), (\sigma_n\otimes\sigma_m)  \right]&= \sum_p2 \,i\varepsilon_{i,n,p}\sigma_p\otimes\delta_{j,m}+\sum_q2\,\delta_{i,n}\otimes\varepsilon_{j,m,q}\sigma_q
\end{align*}
It vanishes if the pairs $i,n$ and $j,m$ consist of different matrices respectively, thus in any configuration, given the restrictions of this case. Applying this consideration to the matrix
\begin{align}\label{eq:momentum Hamiltonian}
    H_\mathbf{k} &=\mathbb{1}\otimes H_\mathbf{k}^{A,A} + \sigma_x\otimes\text{Re}\left[ H_\mathbf{k}^{A,B} \right]-\sigma_y\otimes\text{Im}\left[ H_\mathbf{k}^{A,B} \right]
\end{align}
where the submatrices have the structure
\begin{align*}
    O &= O_0\,\mathbb{1}+ O_x\,\sigma_x+ O_z\,\sigma_z
\end{align*}
Using the fact that all contributions to $H_k$ are independent functions, as stated in the main part, each matrix commuting with all $H_k$ for all $k$ has to commute with each contribution individually.
Starting with the $\mathbb{1}\otimes\sigma_x$ part of the $A,A$ term, the possible contributions to a commuting matrix are restricted to
\begin{align*}
    \sigma_{0/x/y/z}\otimes\sigma_{0/x}
\end{align*}
whereas the part proportional to $\mathbb{1}\otimes\sigma_y$ restricts the possible contributions to 
\begin{align*}
    \sigma_{0/x/y/z}\otimes\sigma_{0}
\end{align*}
Turning to the second term of \eqref{eq:momentum Hamiltonian}, the $\sigma_x\otimes\sigma_{0/x/z}$ part rules out all contributions but
\begin{align*}
    \sigma_{0/x}\otimes\sigma_{0}
\end{align*}
And the last term of \eqref{eq:momentum Hamiltonian} finally only permits multiples of the identity as commuting matrices. Therefore the momentum space Hamiltonian matrix is irreducible.\\\\
Another useful identity for the evaluation of chiral symmetry is
\begin{align*}
    \{\sigma_i\otimes\sigma_j,\sigma_n\otimes\sigma_m\}=\half\,\{\sigma_i,\sigma_n\}\otimes\{\sigma_j,\sigma_m\}
\end{align*}
This can be shown by distinguishing the three exclusive cases from above:
\begin{itemize}
    \item case 1: at least one of the pauli matrices is the identity. Then without loss of generality if $\sigma_i=\mathbb{1}$
    \begin{align*}
        \{\sigma_i\otimes\sigma_j,\sigma_n\otimes\sigma_m\}&=\sigma_n\otimes\sigma_j\sigma_m+\sigma_n\otimes\sigma_m\sigma_j\\
        &=\sigma_n\otimes\{\sigma_j,\sigma_m\}\\
        &=\half\,\{\sigma_i,\sigma_n\}\otimes\{\sigma_j,\sigma_m\}
    \end{align*}
    \item case 2: $\sigma_i=\sigma_n$ or $\sigma_j=\sigma_m$. Assuming the former yields
    \begin{align*}
        \{\sigma_i\otimes\sigma_j,\sigma_n\otimes\sigma_m\}&=\mathbb{1}\otimes\sigma_j\sigma_m+\mathbb{1}\otimes\sigma_m\sigma_j\\
        &=\half\,\{\sigma_i,\sigma_n\}\otimes\{\sigma_j,\sigma_m\}
    \end{align*}
    \item case 3: neither of the previous cases. Here it holds
    \begin{align*}
        \{\sigma_i\otimes\sigma_j,\sigma_n\otimes\sigma_m\}&=\sigma_i\sigma_n\otimes\sigma_j\sigma_m+\sigma_n\sigma_i\otimes\sigma_m\sigma_j\\
        &=\sum_{p,q} i\varepsilon_{inp}\,\sigma_p\otimes i\varepsilon_{jmq}\,\sigma_q+\sum_{p,q} i\varepsilon_{nip}\,\sigma_p\otimes i\varepsilon_{mjq}\,\sigma_q\\
        &=2\,\sum_{p,q} i\varepsilon_{inp}\,\sigma_p\otimes i\varepsilon_{jmq}\,\sigma_q\\
        &=\half[\sigma_i,\sigma_n]\otimes[\sigma_j,\sigma_m]\neq0
    \end{align*}
\end{itemize}